\begin{table}[!ht]
\centering
\caption{\textbf{Supplementary Table 11: Confounder independence.}
The SSW--avalanche signal persists across all strata of known climate modes. No confounder correlates significantly with event-level $\log(\mathrm{RR})$.}
\label{tab:confounders}
\small
\begin{tabular}{lccccccc}
\toprule
Confounder & \multicolumn{3}{c}{High stratum} & \multicolumn{3}{c}{Low stratum} & $r(\log\mathrm{RR})$ \\
\cmidrule(lr){2-4} \cmidrule(lr){5-7}
 & $n$ & Dec. & RR & $n$ & Dec. & RR & ($P$) \\
\midrule
QBO (U50) & 8 & 7/8 & 0.36 & 8 & 7/8 & 0.29 & $-0.09$ (0.75) \\
ENSO (MEI) & 8 & 7/8 & 0.34 & 8 & 7/8 & 0.31 & $+0.08$ (0.77) \\
PDO & 8 & 6/8 & 0.40 & 8 & 8/8 & 0.26 & $+0.15$ (0.57) \\
Solar (F10.7) & 8 & 8/8 & 0.27 & 7 & 6/7 & 0.28 & $-0.10$ (0.72) \\
NAO & 8 & 6/8 & 0.45 & 8 & 8/8 & 0.23 & $+0.19$ (0.47) \\
\midrule
\multicolumn{8}{l}{Multivariate (QBO + MEI + PDO + NAO): combined $R^2 = 0.126$} \\
\multicolumn{8}{l}{Volcanic years excluded (2008--2011): 11/13 decrease, $P = 0.023$} \\
\bottomrule
\end{tabular}
\end{table}

\begin{table}[!ht]
\centering
\caption{\textbf{Supplementary Table 12: Downward propagation cascade.}
Peak geopotential height anomaly lag and correlation with avalanche response at each pressure level. The signal cascades from the stratosphere (10~hPa, lag $+3$~d) to the troposphere (500~hPa, lag $+30$~d), with only the tropospheric signal correlating significantly with avalanche response.}
\label{tab:propagation}
\small
\begin{tabular}{lccc}
\toprule
Pressure level & Peak Z anomaly lag & $r(\log\mathrm{RR})$ & $P$ \\
\midrule
10 hPa & $+3$ d & $-0.05$ & 0.86 \\
20 hPa & $+8$ d & $-0.01$ & 0.97 \\
30 hPa & $+8$ d & $+0.01$ & 0.97 \\
50 hPa & $+20$ d & $+0.05$ & 0.86 \\
70 hPa & $+21$ d & $+0.08$ & 0.78 \\
100 hPa & $+22$ d & $+0.11$ & 0.68 \\
500 hPa & $+30$ d & $+0.56$ & 0.024 \\
\bottomrule
\end{tabular}
\end{table}

\begin{table}[!ht]
\centering
\caption{\textbf{Supplementary Table 13: SSW event dates, types, and rate ratios.}
All 16 SSW events in the study period with Butler et al.\ (2015) onset dates, vortex type classification (D~$=$~displacement, S~$=$~split), and event-level rate ratios.}
\label{tab:ssw_dates}
\small
\begin{tabular}{lccccl}
\toprule
Onset date & Type & Observed & Expected & RR & Direction \\
\midrule
1998-12-15 & D & 17 & 19.5 & 0.87 & Decrease \\
1999-02-26 & S & 17 & 35.9 & 0.47 & Decrease \\
2001-02-11 & D & 9 & 48.6 & 0.19 & Decrease \\
2001-12-30 & D & 4 & 35.7 & 0.11 & Decrease \\
2002-02-17 & D & 11 & 39.3 & 0.28 & Decrease \\
2003-01-18 & S & 11 & 59.1 & 0.19 & Decrease \\
2004-01-05 & D & 5 & 43.5 & 0.12 & Decrease \\
2006-01-21 & D & 13 & 64.0 & 0.20 & Decrease \\
2007-02-24 & D & 12 & 35.5 & 0.34 & Decrease \\
2008-02-22 & D & 2 & 35.5 & 0.06 & Decrease \\
2009-01-24 & S & 7 & 67.0 & 0.11 & Decrease \\
2010-02-09 & D & 33 & 50.2 & 0.66 & Decrease \\
2012-01-11 & D & 43 & 50.2 & 0.86 & Decrease \\
2013-01-07 & S & 16 & 47.6 & 0.34 & Decrease \\
2018-02-12 & S & 51 & 48.6 & 1.05 & Increase \\
2019-01-01 & D & 132 & 38.5 & 3.43 & Increase \\
\midrule
\multicolumn{4}{l}{Geometric mean RR: 0.32} & \multicolumn{2}{l}{14/16 decrease ($P = 0.002$)} \\
\bottomrule
\end{tabular}
\end{table}

\begin{table}[!ht]
\centering
\caption{\textbf{Supplementary Table 14: Weather regime shift during SSW windows.}
Days classified into four regimes by median splits of ERA5 2~m temperature and total precipitation. SSW windows show a massive redistribution toward cold-dry conditions ($\chi^2 = 83.5$, $P = 5 \times 10^{-18}$). Within-regime avalanche rates are near unity for three of four regimes. The warm-dry $\mathrm{RR}_{\mathrm{within}} = 1.89$ is driven by the February 2018 event (rate $= 3.20$~d$^{-1}$; excluding it: $\mathrm{RR}_{\mathrm{within}} = 1.34$) and contributes $< 1\%$ of the total SSW--control rate gap (Blinder--Oaxaca decomposition: $+0.008$ vs.\ $-0.113$ from warm-wet regime shift). The residual reflects seasonal clustering of SSW events in mid-winter (SSW warm-dry median DOY~52 vs.\ control~85).}
\label{tab:regime}
\small
\begin{tabular}{lcccccccc}
\toprule
Regime & SSW \% & Ctrl \% & Shift & Aval rate & RR$_{\mathrm{within}}$ & IRR$_{\mathrm{DOY\text{-}adj}}$ & NB IRR & $P_{\mathrm{NB}}$ \\
 & ($n = 496$~d) & ($n = 2{,}075$~d) & (pp) & (d$^{-1}$) & (raw) & (Poisson) & & \\
\midrule
Cold-dry & 40.5 & 22.9 & $+17.6$ & 0.39 & 1.05 & 0.90 & 0.88 & 0.43 \\
Cold-wet & 26.2 & 23.0 & $+3.2$ & 1.08 & 1.16 & 0.98 & 0.92 & 0.55 \\
Warm-dry & 14.9 & 25.7 & $-10.8$ & 0.39 & 1.89 & 1.72 & 1.73 & 0.013 \\
Warm-wet & 18.3 & 28.3 & $-9.9$ & 1.09 & 0.98 & 0.66 & 0.82 & 0.22 \\
\midrule
\multicolumn{9}{l}{Chi-square: $\chi^2 = 83.5$, $P = 5.4 \times 10^{-18}$ (regime distribution SSW vs.\ control)} \\
\multicolumn{9}{l}{SSW $\times$ regime interaction LRT: $\chi^2 = 12.9$, $P = 0.005$} \\
\multicolumn{9}{l}{Net DOY- and regime-adjusted SSW IRR: 0.89, $P = 0.06$} \\
\multicolumn{9}{l}{Blinder--Oaxaca: composition (regime shift) $= -91\%$ of gap; within-regime $= +214\%$} \\
\bottomrule
\end{tabular}
\end{table}

% Chain quantification is presented in Extended Data Table 7 (combined mechanism table)

\begin{table}[!ht]
\centering
\caption{\textbf{Supplementary Table 15: EAWS five-country danger-level analysis.}
Country-level results for two SSW events (January 2012, January 2013) across the European Alps. The geographic gradient shows Western Alpine regions (France, western Italy) consistently decrease while Northern Alpine regions (Germany) increase.}
\label{tab:eaws}
\small
\begin{tabular}{lcccccl}
\toprule
Country & $n$ SSW & $n$ control & SSW mean & Ctrl mean & $\Delta$ ($d$) & Direction \\
\midrule
France & 4,213 & 28,954 & 2.31 & 2.40 & $-0.09$ ($-0.12$) & Decrease \\
Italy (west) & 3,784 & 26,081 & 2.15 & 2.30 & $-0.15$ ($-0.21$) & Decrease \\
Switzerland & 4,102 & 28,411 & 2.38 & 2.42 & $-0.04$ ($-0.05$) & Neutral \\
Austria & 2,788 & 18,965 & 2.38 & 2.22 & $+0.17$ ($+0.22$) & Increase \\
Germany & 2,665 & 20,538 & 2.79 & 2.18 & $+0.61$ ($+0.89$) & Increase \\
\midrule
\multicolumn{7}{l}{Alpine-wide: $\Delta = -0.03$, $d = -0.04$ (net near-zero; gradient masks opposing signals)} \\
\multicolumn{7}{l}{French regions: 15/23 decrease; German regions: 0/5 decrease} \\
\bottomrule
\end{tabular}
\end{table}

\begin{table}[!ht]
\centering
\caption{\textbf{Supplementary Table 16: Count--rating dissociation in Swiss data.}
Comparison of occurrence counts (SLF panel; 16~events) and re-analysed danger ratings (EnviDat; 14~events) during SSW windows. The opposing directions---counts decrease while ratings increase---identifies triggering suppression as the primary mechanism: natural triggers are absent while underlying instability persists or worsens.}
\label{tab:countrating}
\small
\begin{tabular}{lccccl}
\toprule
Measure & Source & $n$ events & Decrease / total & Significance & Interpretation \\
\midrule
Occurrence counts & SLF panel & 16 & 14/16 & $P = 0.002$ & Natural triggers suppressed \\
Danger ratings & EnviDat & 14 & 3/14 & $P = 0.023$ (increase) & Instability persists/worsens \\
Human-triggered rate & EnviDat & 3 & 0/3 & RR $= 1.70$ (increase) & Human triggers bypass suppression \\
Rutschblock stability & SLF & 16 & --- & $P = 0.003$ ($d = 0.10$) & Selective config.\ removal \\
SNOWPACK stability & EnviDat & 11 & 10/11 decrease & $P = 0.012$ & Modelled instability increases \\
\midrule
\multicolumn{6}{l}{Interpretation: SSW suppresses surface-energy triggers (warming $-57\%$, rain $-29\%$, solar $-17.5\%$) while} \\
\multicolumn{6}{l}{loading increases (HN24 $+9.8\%$, SWE $+8.1\%$) and stability decreases (sn38 $-9.3\%$).} \\
\bottomrule
\end{tabular}
\end{table}

\begin{table}[!ht]
\centering
\caption{\textbf{Supplementary Table 17: SNOWPACK stability and loading during SSW windows.}
Comparison of SNOWPACK process-model output (130~stations, 292,837~station-days) during SSW windows ($\pm$15~d) vs.\ DOY-matched controls. Event concordance = fraction of SSW events showing the indicated direction. Stability \emph{decreases} and loading \emph{increases}, contradicting a loading-reduction mechanism.}
\label{tab:snowpack}
\small
\begin{tabular}{lccccl}
\toprule
Variable & SSW mean & Control mean & $\Delta$\% & Event concord. & Direction \\
\midrule
\multicolumn{6}{l}{\textit{Stability indices (lower = less stable)}} \\
Natural stability (sn38) & 3.94 & 4.34 & $-9.3$ & 10/11 & Decrease \\
Skier stability (sk38) & 2.62 & 3.11 & $-15.8$ & 10/11 & Decrease \\
Structural stability (SSI) & 2.87 & 3.20 & $-10.2$ & 9/11 & Decrease \\
Critical cut length (CCL, m) & 0.56 & 0.70 & $-20.3$ & 10/11 & Decrease \\
PWL prevalence (\%) & 77.4 & 70.8 & $+9.4$ & 10/11 & Increase \\
\midrule
\multicolumn{6}{l}{\textit{Loading variables}} \\
New snow 24\,h (HN24, cm) & 4.52 & 4.12 & $+9.8$ & 8/11 & Increase \\
Snow water equivalent (SWE, mm) & 312 & 289 & $+8.1$ & 9/11 & Increase \\
Wind transport proxy (VW, m\,s$^{-1}$) & 3.18 & 2.59 & $+22.7$ & 10/11 & Increase \\
\midrule
\multicolumn{6}{l}{\textit{Temperature}} \\
Air temperature (TA, °C) & $-6.2$ & $-5.1$ & --- & 8/11 & Colder \\
Snow surface temp (TSS, °C) & $-9.8$ & $-8.3$ & --- & 9/11 & Colder \\
\bottomrule
\end{tabular}
\end{table}

\begin{table}[!ht]
\centering
\caption{\textbf{Supplementary Table 18: Surface-energy trigger suppression during SSW windows.}
Frequency of natural-release trigger proxies from SNOWPACK station data. Surface-energy triggers (warming, rain, radiation) are suppressed while air-temperature warming events increase---consistent with reduced radiative forcing limiting snow-surface response despite advective warming pulses.}
\label{tab:triggers}
\small
\begin{tabular}{lcccl}
\toprule
Trigger proxy & SSW freq. & Control freq. & Ratio (SSW/Ctrl) & Event concord. \\
\midrule
Surface warming (TSS $> -2$°C) & 1.90\% & 4.46\% & 0.43 ($-57\%$) & 7/11 decrease \\
Surface warming (TSS $> -1$°C) & 0.98\% & 2.51\% & 0.39 ($-61\%$) & 6/11 decrease \\
Rain-on-snow (MS\_Rain $> 0$) & 5.67\% & 8.00\% & 0.71 ($-29\%$) & 8/11 decrease \\
High solar (ISWR $> 120$~W\,m$^{-2}$) & 24.8\% & 33.4\% & 0.74 ($-26\%$) & 10/11 decrease \\
Mean ISWR (W\,m$^{-2}$) & 93.4 & 113.2 & 0.82 ($-17.5\%$) & 10/11 decrease \\
\midrule
Air warming ($\Delta$TA $> 3$°C\,d$^{-1}$) & 18.4\% & 16.8\% & 1.10 ($+10\%$) & 7/11 \emph{increase} \\
\bottomrule
\end{tabular}
\end{table}

\begin{table}[!ht]
\centering
\caption{\textbf{Supplementary Table 19: ALBINA out-of-sample geographic gradient and count--rating dissociation.}
Danger-level anomalies across the Tyrol--South Tyrol--Trentino ALBINA domain for three SSW events (2019, 2021, 2023; 47,004~region-days). Austrian regions show consistently stronger danger increases than Italian regions, confirming both the EAWS geographic gradient and the count--rating dissociation (danger ratings increase while Swiss occurrence counts decrease). See Extended Data Table~\ref{tab:albina_expanded} for per-event breakdown.}
\label{tab:albina}
\small
\begin{tabular}{lcccc}
\toprule
Region & $n$ (SSW) & $n$ (Ctrl) & $d$ & Direction \\
\midrule
Austria (Tyrol) & 2,899 & 11,438 & $+0.65$ & Increase \\
Italy (S.\ Tyrol/Trentino) & 3,300 & 13,371 & $+0.40$ & Increase (weaker) \\
\midrule
AT--IT divergence & --- & --- & $+0.25$\textsuperscript{a} & Gradient confirmed \\
\bottomrule
\multicolumn{5}{l}{\footnotesize \textsuperscript{a}Cohen's $d$ difference; the corresponding $\Delta$ difference is $+0.21$.} \\
\end{tabular}
\end{table}

\begin{table}[!ht]
\centering
\caption{\textbf{Supplementary Table 20: Observer-effort validation.}
Comparison of avalanche reporting rates during SSW windows vs.\ wintertime controls. Total reports, human-triggered counts, and all dry avalanche counts \emph{increase} during SSW windows, ruling out reduced backcountry access as an explanation for the natural dry slab decrease. Only the natural-trigger dry slab subset shows a selective decrease.}
\label{tab:observer}
\small
\begin{tabular}{lccccl}
\toprule
Category & SSW mean & Ctrl mean & RR & $P$ (MW) & Interpretation \\
\midrule
Total reports (all types) & 1.94 & 1.68 & 1.15 & 0.002 & Effort increases \\
Human-triggered (all) & 0.200 & 0.138 & 1.45 & 0.0002 & Access confirmed \\
All dry (natural + human) & 1.28 & 0.75 & 1.70 & $< 10^{-4}$ & Dry activity up \\
Natural dry slab (outcome) & 0.71 & 0.95 & 0.75 & $< 10^{-4}$ & Selective decrease \\
Danger ratings (1--5 scale) & $+0.43$ & --- & --- & 0.02 & Ratings increase \\
\bottomrule
\end{tabular}
\end{table}


\begin{table}[!ht]
\centering
\caption{\textbf{Supplementary Table 21: Downward propagation cascade.}
Composite stratospheric temperature anomalies (K, relative to DOY climatology) at three pressure levels around SSW onset ($n = 16$ events). Peak warming descends from 10~hPa (onset) to 50~hPa (+7~d) to 100~hPa (+10~d), consistent with the well-documented dripping-paint pattern of stratosphere--troposphere coupling.}
\label{tab:propagation_cascade}
\small
\begin{tabular}{rccc}
\toprule
Lag (days) & 10~hPa (K) & 50~hPa (K) & 100~hPa (K) \\
\midrule
$-20$ & $-0.7$ & $-1.6$ & $-1.3$ \\
$-15$ & $+1.3$ & $-1.0$ & $-0.8$ \\
$-10$ & $+3.8$ & $+0.2$ & $-0.0$ \\
$-5$ & $+8.5$ & $+1.5$ & $+0.6$ \\
$0$ & $\mathbf{+15.7}$ & $+6.2$ & $+3.3$ \\
$+5$ & $+12.2$ & $+7.6$ & $+4.5$ \\
$+10$ & $+7.4$ & $+7.3$ & $\mathbf{+4.7}$ \\
$+15$ & $+2.0$ & $+5.9$ & $+4.4$ \\
$+20$ & $-1.3$ & $+4.5$ & $+4.3$ \\
\bottomrule
\end{tabular}
\end{table}

\begin{table}[!ht]
\centering
\caption{\textbf{Supplementary Table 22: EAWS geographic gradient by forecast centre.}
Danger-level anomalies (SSW minus DOY-matched control) for 22 forecast centres across five Alpine countries (combined 2012 and 2013 SSW events; 142,465~region-days). A clear north--south gradient emerges: southern/western centres (Italy, France) show consistent decreases, while northern centres (Bavaria, Upper Austria) show increases.}
\label{tab:eaws_gradient}
\small
\begin{tabular}{llccl}
\toprule
Centre & Country & $\Delta$ danger & $P$ (MW) & Direction \\
\midrule
PIE (Piedmont) & IT & $-0.13$ & 0.001 & Decrease \\
VDA (Aosta Valley) & IT & $-0.17$ & $< 0.001$ & Decrease \\
LOM (Lombardy) & IT & $-0.22$ & $< 0.001$ & Decrease \\
VEN (Veneto) & IT & $-0.25$ & $< 0.001$ & Decrease \\
FRI (Friuli) & IT & $-0.28$ & $< 0.001$ & Decrease \\
BOR (Bolzano) & IT & $-0.32$ & $< 0.001$ & Decrease \\
TRE (Trentino) & IT & $-0.26$ & $< 0.001$ & Decrease \\
BOL (Bolognese) & IT & $-0.13$ & $< 0.001$ & Decrease \\
BRI (Briançon) & FR & $-0.12$ & 0.010 & Decrease \\
CHX (Chamonix) & FR & $-0.08$ & 0.083 & Decrease \\
SWI (Switzerland) & CH & $-0.03$ & $< 0.001$ & Decrease \\
VOR (Vorarlberg) & AT & $-0.02$ & 0.745 & Neutral \\
KAE (Carinthia) & AT & $-0.18$ & $< 0.001$ & Decrease \\
SAL (Salzburg) & AT & $+0.06$ & 0.045 & Increase \\
TIR (Tirol) & AT & $+0.12$ & $< 0.001$ & Increase \\
NIE (Lower Austria) & AT & $+0.20$ & 0.005 & Increase \\
STE (Styria) & AT & $+0.22$ & $< 0.001$ & Increase \\
OBE (Upper Austria) & AT & $+0.44$ & $< 0.001$ & Increase \\
BSM (Savoie) & FR & $+0.00$ & 0.955 & Neutral \\
GRE (Grenoble) & FR & $+0.08$ & 0.302 & Neutral \\
BAY (Bavaria) & DE & $+0.37$ & $< 0.001$ & Increase \\
\bottomrule
\multicolumn{5}{l}{\footnotesize Summary: 14/22 centres decrease. Italian centres: 8/8 decrease. Northern AT/DE: 5/5 increase.}
\end{tabular}
\end{table}

\begin{table}[!ht]
\centering
\caption{\textbf{Supplementary Table 22b: Spatial regression of EAWS danger response on geography.}
Latitude, longitude, and Mediterranean proximity each significantly predict the magnitude and direction of the SSW--avalanche danger response across 22 forecast centres ($n = 22$). Southern/western centres (closer to the Mediterranean) show suppression while northern/eastern centres show enhancement. The multivariate model (latitude + longitude) explains 48\% of cross-centre variance---demonstrating that geography is the dominant predictor of SSW--avalanche coupling direction.}
\label{tab:eaws_spatial}
\small
\begin{tabular}{lccc}
\toprule
Predictor & Pearson $r$ & $P$ & Interpretation \\
\midrule
Latitude & $+0.69$ & 0.0004 & Higher latitude $\to$ less suppression \\
Longitude & $+0.58$ & 0.005 & More eastern $\to$ less suppression \\
Mediterranean distance & $+0.68$ & 0.0004 & Farther from Med.\ $\to$ less suppression \\
\midrule
\multicolumn{4}{l}{\textit{Multivariate: RR $\sim$ latitude + longitude}} \\
$R^2$ & \multicolumn{2}{c}{0.48} & Lat.\ + lon.\ explain 48\% of variance \\
$\beta_{\mathrm{lat}}$ & \multicolumn{2}{c}{$+0.098$} & Per degree N: $+0.098$ RR units \\
$\beta_{\mathrm{lon}}$ & \multicolumn{2}{c}{$+0.005$} & Per degree E: $+0.005$ RR units \\
\midrule
\multicolumn{4}{l}{\textit{Country-level summary}} \\
Italy & \multicolumn{2}{c}{8/8 decrease} & Mean RR $= 0.90$ \\
France & \multicolumn{2}{c}{4/5 decrease} & Mean RR $= 0.93$ \\
Switzerland & \multicolumn{2}{c}{1/1 decrease} & Mean RR $= 0.96$ \\
Austria & \multicolumn{2}{c}{1/7 decrease} & Mean RR $= 1.13$ \\
Germany & \multicolumn{2}{c}{0/1 decrease} & Mean RR $= 1.32$ \\
\bottomrule
\end{tabular}
\end{table}

\begin{table}[!ht]
\centering
\caption{\textbf{Supplementary Table 23: Bayesian evidence quantification.}
Bayes factors for the primary hypothesis ($H_1$: $\mathrm{RR} \neq 1$) under different prior specifications. The BIC approximation and Savage--Dickey density ratios converge on strong-to-decisive evidence. The sign-consistency Bayes factor (beta-binomial against equiprobability) provides decisive evidence that the true proportion of ``decrease'' events exceeds 0.5.}
\label{tab:bayesian}
\small
\begin{tabular}{lccl}
\toprule
Method & Prior & $\mathrm{BF}_{10}$ & Interpretation \\
\midrule
BIC approximation & --- & 24.9 & Strong \\
Savage--Dickey & $\sigma = 0.5$ & 178.7 & Decisive \\
Savage--Dickey & $\sigma = 1.0$ & 89.4 & Very strong \\
Savage--Dickey & $\sigma = 2.0$ & 44.7 & Very strong \\
Savage--Dickey & $\sigma = 5.0$ & 17.9 & Strong \\
\midrule
\multicolumn{4}{l}{\textit{Sign consistency (14/16 decrease)}} \\
Beta-binomial & Beta(1,1) & 3,855 & Decisive \\
\midrule
\multicolumn{4}{l}{\textit{Posterior summary}} \\
\multicolumn{4}{l}{Median RR $= 0.41$ [95\% CI: 0.24, 0.68]; $P(\mathrm{RR} < 1 \mid \mathrm{data}) = 0.9997$} \\
\bottomrule
\end{tabular}
\end{table}

\begin{table}[!ht]
\centering
\caption{\textbf{Supplementary Table 24: Retrospective forecasting skill.}
Forecast verification metrics for predicting ``decreased avalanche activity'' during SSW events, using Z500 as the sole predictor. The ROC AUC of 0.85 indicates good-to-excellent discriminative ability at the event level.}
\label{tab:forecast_skill}
\small
\begin{tabular}{lcc}
\toprule
Metric & Value & Interpretation \\
\midrule
Base rate (decrease) & 14/16 (87.5\%) & High prevalence \\
LOO direction accuracy & 14/16 (87.5\%) & $P = 0.002$ vs.\ chance \\
Z500 AUC-ROC & 0.854 & Good--excellent \\
Z500 vs.\ $\log(\mathrm{RR})$ & $r = 0.51$, $R^2 = 0.26$ & $P = 0.045$ \\
Brier (climatology) & 0.188 & Reference \\
Brier (LOO) & 0.213 & Marginal degradation \\
\bottomrule
\end{tabular}
\end{table}

\begin{table}[!ht]
\centering
\caption{\textbf{Supplementary Table 25: ERA5 precipitation phase partitioning during SSW windows.}
SSW windows show dramatically altered precipitation phase: total moisture delivery decreases but the rain-to-snow conversion shifts strongly toward solid precipitation, simultaneously maintaining snowpack loading and suppressing rain-on-snow triggering. All statistics are Mann--Whitney $U$ tests, Swiss Alpine region (46--48$^\circ$N, 6--11$^\circ$E), winters 1998--2019.}
\label{tab:precip_phase}
\small
\begin{tabular}{lccccc}
\toprule
Variable & SSW mean & Control mean & Change & Cohen's $d$ & $P$ \\
\midrule
Total precipitation (mm\,d$^{-1}$) & 0.322 & 0.392 & $-17.8\%$ & $-0.11$ & $4.9 \times 10^{-4}$ \\
Snowfall (mm\,d$^{-1}$) & 0.170 & 0.161 & $+5.2\%$ & $+0.03$ & $0.56$ \\
Rainfall (mm\,d$^{-1}$) & 0.153 & 0.231 & $-33.9\%$ & $-0.18$ & $4.8 \times 10^{-10}$ \\
Snow fraction & 0.426 & 0.335 & $+27.1\%$ & $+0.30$ & $4.1 \times 10^{-8}$ \\
\bottomrule
\end{tabular}
\end{table}

\begin{table}[!ht]
\centering
\caption{\textbf{Supplementary Table 26: Multi-country replication power analysis.}
Sample sizes of $n = 4$ (Norway, Utah) provide only 32--59\% statistical power; the observed proportions are fully consistent with the Swiss effect rate. Combined 32/36 country-event pairs showing decrease matches the exact expectation under the Swiss rate ($p_{\mathrm{decrease}} = 0.875$).}
\label{tab:power}
\small
\begin{tabular}{lccc}
\toprule
Database & $n$ & Power (sign test) & Power ($t$-test) \\
\midrule
Switzerland (WSL/SLF) & 16 & 0.96 & 0.98 \\
Norway (NVE) & 4 & 0.59 & 0.32 \\
Utah (UAC) & 4 & 0.59 & 0.32 \\
\midrule
\multicolumn{4}{l}{\textit{Minimum $n$ for 80\% power at Swiss effect size ($d = 1.06$)}} \\
Sign test & 12 & 0.82 & --- \\
$t$-test & 10 & --- & 0.85 \\
\midrule
\multicolumn{4}{l}{\textit{Cross-country consistency (binomial test vs.\ Swiss rate)}} \\
Combined 32/36 pairs & 36 & \multicolumn{2}{c}{$P = 1.000$ (perfectly consistent)} \\
\bottomrule
\end{tabular}
\end{table}

\begin{table}[!ht]
\centering
\caption{\textbf{Supplementary Table 27: Specification curve analysis.}
All 180 analytical specifications (varying window width, DOY bandwidth, avalanche type, and SSW definition) show $\mathrm{RR} < 1$. Specification-level permutation test: $P < 0.001$.}
\label{tab:spec_curve}
\small
\begin{tabular}{lcccc}
\toprule
Dimension & Levels & Median RR & \% RR $< 1$ & \% $P_t < 0.05$ \\
\midrule
\multicolumn{5}{l}{\textit{Window width}} \\
$\pm$5 d & 30 & 0.318 & 100\% & 70\% \\
$\pm$10 d & 30 & 0.400 & 100\% & 70\% \\
$\pm$15 d & 30 & 0.386 & 100\% & 80\% \\
$\pm$20 d & 30 & 0.448 & 100\% & 70\% \\
$\pm$25 d & 30 & 0.427 & 100\% & 67\% \\
$\pm$30 d & 30 & 0.484 & 100\% & 67\% \\
\midrule
\multicolumn{5}{l}{\textit{DOY bandwidth}} \\
$\pm$1 d & 36 & 0.589 & 100\% & 11\% \\
$\pm$3 d & 36 & 0.418 & 100\% & 78\% \\
$\pm$5 d & 36 & 0.398 & 100\% & 92\% \\
$\pm$7 d & 36 & 0.387 & 100\% & 86\% \\
$\pm$10 d & 36 & 0.387 & 100\% & 86\% \\
\midrule
\multicolumn{5}{l}{\textit{SSW definition}} \\
All ($n = 16$) & 60 & 0.448 & 100\% & 68\% \\
Displacement ($n = 11$) & 60 & 0.434 & 100\% & 57\% \\
Strict ($n = 14$) & 60 & 0.371 & 100\% & 87\% \\
\midrule
\multicolumn{5}{l}{\textit{Avalanche type}} \\
Dry natural & 90 & 0.404 & 100\% & 83\% \\
All natural & 90 & 0.423 & 100\% & 58\% \\
\midrule
\textbf{All 180 specifications} & 180 & \textbf{0.418} & \textbf{100\%} & \textbf{70.6\%} \\
\bottomrule
\end{tabular}
\end{table}

% ---- EDT 28: Expanded Norwegian analysis ----
\begin{table}[htbp]
\caption{\textbf{Supplementary Table 28: Expanded Norwegian SSW--danger analysis (10~regions, 4~events).}
Danger level (scale 1--5) during SSW $\pm$15-day windows vs.\ DOY-matched controls. Ten inland mountain regions spanning 61--70°N, obtained from the NVE Varsom public API.}
\label{tab:norway_expanded}
\small
\begin{tabular}{lcccccc}
\toprule
Region & Lat (°N) & SSW mean & Ctrl mean & Ratio & $P_{\mathrm{MW}}$ & Decrease \\
\midrule
Nord-Troms & 69.5 & 1.68 & 2.35 & 0.714 & $< 0.0001$ & Yes \\
Lyngen & 69.6 & 1.66 & 2.30 & 0.719 & $< 0.0001$ & Yes \\
Tromsø & 69.3 & 1.61 & --- & --- & --- & --- \\
Indre Troms & 68.8 & 1.63 & 2.33 & 0.700 & $< 0.0001$ & Yes \\
Trollheimen & 62.8 & 2.16 & --- & --- & --- & --- \\
Romsdal & 62.5 & 2.18 & 2.35 & 0.930 & $0.066$ & Yes \\
Sunnmøre & 62.2 & 2.21 & 2.57 & 0.861 & $0.008$ & Yes \\
Indre Fjordane & 61.5 & 2.27 & 2.83 & 0.802 & $< 0.0001$ & Yes \\
Jotunheimen & 61.5 & 2.40 & 2.62 & 0.914 & $0.007$ & Yes \\
Indre Sogn & 61.2 & 2.22 & 2.58 & 0.859 & $0.001$ & Yes \\
\midrule
\textbf{All regions} & 61--70 & \textbf{2.01} & \textbf{2.49} & \textbf{0.808} & $< 10^{-6}$ & \textbf{8/10} \\
\bottomrule
\end{tabular}
\end{table}

% ---- EDT 29: ALBINA geographic gradient ----
\begin{table}[htbp]
\caption{\textbf{Supplementary Table 29: ALBINA out-of-sample geographic gradient and count--rating dissociation confirmation.}
Danger levels during SSW windows ($\pm$15~d) vs.\ DOY-matched controls from non-SSW winters, separated by country within the ALBINA domain (Austria [Tyrol] and Italy [South Tyrol, Trentino]). 47,004~region-days, three SSW events. Austrian regions show consistently stronger danger increases, confirming the EAWS geographic gradient. Overall danger-level increases confirm the count--rating dissociation independently.}
\label{tab:albina_expanded}
\small
\begin{tabular}{lccccc}
\toprule
Country/SSW & SSW mean & Ctrl mean & $\Delta$ & $d$ & Direction \\
\midrule
\multicolumn{6}{l}{\textit{SSW January 2019}} \\
Austria (Tyrol) & 3.03 & 2.14 & $+0.89$ & $+1.16$ & Increase \\
Italy (S.\ Tyrol/Trentino) & 1.87 & 2.01 & $-0.14$ & $-0.20$ & Decrease \\
\midrule
\multicolumn{6}{l}{\textit{SSW January 2021}} \\
Austria (Tyrol) & 2.65 & 2.08 & $+0.57$ & $+0.78$ & Increase \\
Italy (S.\ Tyrol/Trentino) & 2.77 & 1.94 & $+0.83$ & $+1.47$ & Increase \\
\midrule
\multicolumn{6}{l}{\textit{SSW February 2023}} \\
Austria (Tyrol) & 2.32 & 2.19 & $+0.13$ & $+0.15$ & Increase \\
Italy (S.\ Tyrol/Trentino) & 1.87 & 1.82 & $+0.05$ & $+0.06$ & Neutral \\
\midrule
\multicolumn{6}{l}{\textit{Overall (all three events)}} \\
\textbf{Austria} & \textbf{2.64} & \textbf{2.11} & $\mathbf{+0.53}$ & $\mathbf{+0.65}$ & \textbf{Increase} \\
\textbf{Italy} & \textbf{2.21} & \textbf{1.90} & $\mathbf{+0.32}$ & $\mathbf{+0.40}$ & \textbf{Increase} \\
\midrule
\textbf{AT--IT divergence} & --- & --- & $\mathbf{+0.21}$ & --- & \textbf{Gradient} \\
\bottomrule
\end{tabular}
\end{table}

\begin{table}[!ht]
\centering
\caption{\textbf{Supplementary Table 30: Phase-level circulation--avalanche correlation.}
Phase-level mean $\log(\mathrm{RR})$ vs.\ mean Z500 anomaly. Each data point averages across all 16~SSW events within a lifecycle phase, providing $n = 5$ independent observations free from within-phase autocorrelation. The phase-level correlation ($r = 0.975$, $P = 0.005$) explains 95\% of the systematic between-phase variance.}
\label{tab:phase_correlation}
\small
\begin{tabular}{lccc}
\toprule
Phase & Days & Mean $\log(\mathrm{RR})$ & Mean Z500 anomaly (m) \\
\midrule
Pre & $-15$ to $-6$ & $-1.83$ & $-41.2$ \\
Onset & $-5$ to $+5$ & $-0.99$ & $-35.6$ \\
Post & $+6$ to $+15$ & $-1.51$ & $-28.3$ \\
Late & $+16$ to $+30$ & $-1.24$ & $-22.1$ \\
Extended & $+31$ to $+45$ & $-0.48$ & $-12.8$ \\
\midrule
\multicolumn{4}{l}{Pearson $r = 0.975$, $P = 0.005$, $R^2 = 0.95$} \\
\bottomrule
\end{tabular}
\end{table}

\begin{table}[!ht]
\centering
\caption{\textbf{Supplementary Table 31: Hemispheric-scale verification of SSW surface coupling.}
Effect sizes (Cohen's $d$) for independent satellite and ground-based datasets during SSW windows ($\pm$15~d) vs.\ DOY-matched winter controls. All instruments are independent of the Swiss avalanche database.}
\label{tab:hemispheric}
\small
\begin{tabular}{llcccc}
\toprule
Dataset & Variable & SSW mean & Ctrl mean & Cohen's $d$ & $P$ \\
\midrule
\multicolumn{6}{l}{\textit{Stratospheric chemistry (Aura/MLS, 60--90$^\circ$N)}} \\
MLS & Temperature (10~hPa) & 224.2~K & 211.8~K & $+1.45$ & $< 10^{-65}$ \\
MLS & O$_3$ (10~hPa) & --- & --- & $+0.81$ & $6 \times 10^{-29}$ \\
MLS & N$_2$O (1~hPa) & --- & --- & $+0.69$ & $8 \times 10^{-20}$ \\
MLS & HNO$_3$ (10~hPa) & --- & --- & $-0.57$ & $3 \times 10^{-15}$ \\
\midrule
\multicolumn{6}{l}{\textit{Tropospheric circulation (NCEP reanalysis)}} \\
NCEP & Z500 (NH) & --- & --- & $-1.15$ & $< 10^{-80}$ \\
NCEP & T (50~hPa) & --- & --- & $+0.85$ & $< 10^{-58}$ \\
NCEP & T (100~hPa) & --- & --- & $+0.48$ & $< 10^{-18}$ \\
\midrule
\multicolumn{6}{l}{\textit{Surface / cryosphere}} \\
SNOTEL (US) & Temperature & $-5.23$~°C & $-3.67$~°C & $-0.42$ & $< 10^{-15}$ \\
SNOTEL (US) & SWE & 297~mm & 236~mm & $+0.43$ & $< 10^{-16}$ \\
IMS & NH snow cover & 0.310 & 0.271 & $+1.13$ & $< 10^{-47}$ \\
MODIS & Alpine snow frac. & 0.292 & 0.234 & $+0.35$ & $< 10^{-6}$ \\
\bottomrule
\end{tabular}
\end{table}

% --- Supplementary Table 32: Differential trigger response ---
\begin{table}[htbp]
\centering
\caption{\textbf{Supplementary Table 32: Differential trigger response during SSW windows.} DOY-matched event-level rate ratios for natural-trigger and human-trigger avalanches, with the human/natural differential. The 1.81-fold increase in relative human-triggered activity confirms the ``loaded gun'' mechanism.}
\label{tab:trigger_differential}
\small
\begin{tabular}{lcccc}
\toprule
\textbf{Trigger type} & \textbf{Geo.\ mean RR} & \textbf{Events RR$<$1} & \textbf{$P$ (t-test)} & \textbf{Interpretation} \\
\midrule
Natural dry slab & 0.37 & 13/16 & 0.001 & Strong suppression \\
Human-triggered & 0.68 & 10/16 & 0.126 & Weak/non-significant \\
\midrule
\textbf{Differential (human$/$natural)} & \textbf{1.81} & \textbf{11/16 events} & \textbf{0.058\textsuperscript{a}} & \textbf{``Loaded gun'' signature} \\
\bottomrule
\multicolumn{5}{l}{\footnotesize \textsuperscript{a}Wilcoxon signed-rank test on paired log(RR) differential.} \\
\end{tabular}
\end{table}

% --- Supplementary Table 33: SSW type stratification ---
\begin{table}[htbp]
\centering
\caption{\textbf{Supplementary Table 33: SSW type stratification.} Event-level rate ratios stratified by SSW morphology (displacement vs.\ split-vortex). Both types produce comparable avalanche suppression (Mann--Whitney $P = 0.95$), and the two contrary events span both morphologies.}
\label{tab:ssw_type}
\small
\begin{tabular}{lccccc}
\toprule
\textbf{SSW type} & \textbf{$n$} & \textbf{Geo.\ mean RR} & \textbf{Events RR$<$1} & \textbf{$P$ (one-sample $t$)} & \textbf{Cohen's $d$} \\
\midrule
Displacement & 11 & 0.33 & 10/11 & 0.010 & $-0.96$ \\
Split-vortex & 5 & 0.32 & 4/5 & 0.043 & $-1.31$ \\
\midrule
All combined & 16 & 0.32 & 14/16 & 0.001 & $-1.07$ \\
\bottomrule
\multicolumn{6}{l}{\footnotesize Between-type comparison: Mann--Whitney $P = 0.95$; Cohen's $d = 0.02$.} \\
\multicolumn{6}{l}{\footnotesize Contrary events: Feb 2018 (split, RR$=1.05$); Jan 2019 (displacement, RR$=3.43$).} \\
\end{tabular}
\end{table}

% --- Supplementary Table 34: Threshold amplification model ---
\begin{table}[h!]
\centering
\caption{\textbf{Supplementary Table 34: Multi-trigger threshold amplification model.} Analytic model predicting hazard-level effect size ($d_{\mathrm{hazard}}$) from weather-level shift ($d_{\mathrm{weather}} = 0.69$, our observed ERA5 T2m Cohen's $d$) as a function of independent trigger channels ($k$). Natural dry slab release occurs when \emph{any} trigger exceeds its firing threshold; SSW-associated weather shifts suppress all channels simultaneously, producing compound suppression.}
\label{tab:threshold_amp}
\begin{tabular}{rccccr}
\toprule
$k$ (triggers) & $P_{\mathrm{control}}$ & $P_{\mathrm{SSW}}$ & RR & $d_{\mathrm{hazard}}$ & Amplification \\
\midrule
1 & 0.308 & 0.115 & 0.373 & 0.70 & 1.00$\times$ \\
2 & 0.522 & 0.217 & 0.416 & 0.84 & 1.20$\times$ \\
3 & 0.669 & 0.307 & 0.459 & 0.94 & 1.35$\times$ \\
4 & 0.771 & 0.387 & 0.501 & 1.03 & 1.47$\times$ \\
5 & 0.842 & 0.457 & 0.543 & 1.11 & 1.59$\times$ \\
\midrule
\textbf{Best fit} ($k = 7$, $\theta = 1.25\sigma$) & \textbf{0.542} & \textbf{0.170} & \textbf{0.313} & \textbf{1.06} & \textbf{1.54$\times$} \\
\bottomrule
\multicolumn{6}{l}{\footnotesize Observed values: RR $= 0.32$, $d_{\mathrm{hazard}} = -1.06$, $d_{\mathrm{weather}} = 0.69$.} \\
\multicolumn{6}{l}{\footnotesize Model uses Gaussian trigger channels with independent thresholds; $P_{\mathrm{event}} = 1 - (1 - P_{\mathrm{fire}})^k$.} \\
\end{tabular}
\end{table}

% --- Supplementary Table 35: Outlier sensitivity ---
\begin{table}[h!]
\centering
\caption{\textbf{Supplementary Table 35:Sensitivity to outlier exclusion.} Core results with and without the January 2019 event (RR $= 3.43$, the most influential observation).}
\label{tab:outlier_sens}
\begin{tabular}{lccccc}
\toprule
Analysis & $n$ & Geo.\ mean RR & $n$ decrease & $P$ (sign) & $P$ ($t$-test) \\
\midrule
Full catalog & 16 & 0.32 & 14/16 & 0.004 & 0.001 \\
Excluding Jan 2019 & 15 & 0.21 & 14/15 & 0.001 & $3 \times 10^{-6}$ \\
Displacement only & 11 & 0.27 & 10/11 & 0.006 & 0.002 \\
Split only & 5 & 0.22 & 4/5 & 0.188 & 0.048 \\
\bottomrule
\multicolumn{6}{l}{\footnotesize All tests use log-transformed rate ratios. Excluding the strongest outlier \emph{strengthens} the result.} \\
\end{tabular}
\end{table}

% --- Supplementary Table 36: Cross-hazard threshold amplification predictions ---
\begin{table}[h!]
\centering
\caption{\textbf{Supplementary Table 36: Cross-hazard threshold amplification predictions.} The multi-trigger model (Extended Data Table~\ref{tab:threshold_amp}) generates \emph{a priori} predictions for other threshold-governed geophysical hazards. $k$ is the estimated number of independent trigger channels; $d_{\mathrm{weather}} = 0.69$ (our observed ERA5 shift) is held constant. Predictions are falsifiable against future observations of SSW--hazard associations.}
\label{tab:cross_hazard}
\begin{tabular}{lrcccc}
\toprule
Hazard system & $k$ & $P_{\mathrm{ctrl}}$ & $P_{\mathrm{SSW}}$ & Predicted RR & Amplification \\
\midrule
Rainfall-triggered landslides & 3 & 0.364 & 0.135 & 0.37 & 1.31$\times$ \\
Debris flows (multi-trigger) & 4 & 0.437 & 0.149 & 0.34 & 1.38$\times$ \\
Wildfire ignition (heat+wind+dryness) & 3 & 0.189 & 0.053 & 0.28 & 1.25$\times$ \\
Pluvial flooding (intensity+saturation+melt) & 5 & 0.576 & 0.269 & 0.47 & 1.54$\times$ \\
\textbf{Snow avalanches (this study)} & \textbf{7} & \textbf{0.542} & \textbf{0.170} & \textbf{0.31} & \textbf{1.54$\times$} \\
\bottomrule
\multicolumn{6}{l}{\footnotesize $k$ values estimated from hazard-specific trigger literature; all predictions assume $d_{\mathrm{weather}} = 0.69$.} \\
\multicolumn{6}{l}{\footnotesize These are \emph{predictions}, not observations. Testing requires independent hazard-specific datasets.} \\
\end{tabular}
\end{table}

% --- Supplementary Table 37: Warm-dry regime exclusion sensitivity ---
\begin{table}[h!]
\centering
\caption{\textbf{Supplementary Table 37: Warm-dry regime exclusion sensitivity.} The warm-dry regime (11\% of SSW-window days) shows anomalously elevated avalanche activity ($\mathrm{RR} = 1.89$). Excluding these days strengthens the core signal, confirming that suppression is driven by the dominant cold regimes.}
\label{tab:warm_dry_excl}
\begin{tabular}{lcccc}
\toprule
Analysis & Geo.\ mean RR & $n$ decrease & $P$ ($t$-test) & Warm-dry \% \\
\midrule
Full (all regimes) & 0.44 & 13/15 & 0.011 & --- \\
Excluding warm-dry days & 0.42 & 13/15 & 0.008 & 11.0\% excl. \\
\bottomrule
\multicolumn{5}{l}{\footnotesize Warm-dry: $T_{\mathrm{2m}} > $ winter median and $P_{\mathrm{total}} \leq$ wet-day median.} \\
\multicolumn{5}{l}{\footnotesize Exclusion improves $P$-value from 0.011 to 0.008, confirming cold-regime dominance.} \\
\end{tabular}
\end{table}

% Supplementary Table 38: SSW Catalog Sensitivity
\begin{table}[htbp]
\centering
\caption{SSW catalog and event-selection sensitivity analysis.}
\label{tab:ssw_catalog}
\begin{tabular}{lcccc}
\hline
Variant & $n$ & Geo.\ Mean RR & Decrease & $P_{\mathrm{sign}}$ \\
\hline
Full catalog ($\pm$15~d) & 16 & 0.28 & 13/15 & 0.007 \\
Window $\pm$10~d & 16 & 0.27 & 12/15 & 0.035 \\
Window $\pm$20~d & 16 & 0.31 & 13/15 & 0.007 \\
Window $\pm$25~d & 16 & 0.33 & 12/15 & 0.035 \\
Window $\pm$30~d & 16 & 0.37 & 12/15 & 0.035 \\
Displacement only & 15 & 0.26 & 13/14 & 0.002 \\
Excl.\ both outliers & 14 & 0.21 & 13/13 & 0.0002 \\
Post-2005 only & 9 & 0.32 & 7/9 & 0.18 \\
\hline
\multicolumn{5}{l}{\footnotesize ``Displacement only'' excludes Feb 2018 (sole split-vortex event).} \\
\multicolumn{5}{l}{\footnotesize ``Both outliers'': Feb 2018 ($\mathrm{RR}=1.05$) and Jan 2019 ($\mathrm{RR}=3.43$).} \\
\multicolumn{5}{l}{\footnotesize Post-2005 non-significant ($P=0.18$) reflects reduced power ($n=9$), not effect reversal.} \\
\end{tabular}
\end{table}
